\chapter{Schwarz–Christoffel Transformations}
\label{ch:iv18}

Schwarz–Christoffel transformations map canonical half-plane or disk domains to polygonal regions. The exponents in the derivative encode turning angles at polygon vertices.

\section*{Learning goals}
After this chapter, the reader should be able to:
\begin{itemize}
\item derive Schwarz--Christoffel exponents from polygon interior angles;
\item construct the derivative of a Schwarz--Christoffel map from prevertices;
\item identify accessory constants and normalization conditions in polygon-mapping problems;
\item analyze local behavior near polygon vertices from the power exponents;
\item map standard half-plane or disk domains to simple polygons;
\item check orientation and vertex order before interpreting a Schwarz--Christoffel formula;
\end{itemize}
\section*{Conceptual roadmap}
\[
\boxed{\text{local holomorphic structure}\;\longrightarrow\;\text{integral or series representation}\;\longrightarrow\;\text{global consequence}.}
\]

\section{Polygonal target}
A polygon is specified by ordered vertices and interior angles $\alpha_k\pi$.

% BEGIN VOL04-EXPANSION IV18-example-01
\begin{example}[Interior angle determines the exponent]\label{ex:iv18-expand-01}
Near a prevertex \(x_k\), write the polygon interior angle as \(\alpha_k\pi\). A Schwarz--Christoffel derivative then has local form
\[
f'(z)\sim C(z-x_k)^{\alpha_k-1}.
\]
Integrating gives \(f(z)-f(x_k)\sim C'(z-x_k)^{\alpha_k}\), so a half-plane angle \(\pi\) is transformed into the polygon angle \(\alpha_k\pi\). This is the same normalized angle convention used throughout the chapter.
\end{example}
% END VOL04-EXPANSION IV18-example-01
\section{Derivative form}
For a half-plane map with real prevertices $x_k$, the derivative has factors $(z-x_k)^{\alpha_k-1}$.

\section{Turning angles}
The exponent $\alpha_k-1$ records the change of boundary direction at the corresponding vertex.

\section{Integration}
The conformal map is obtained by integrating the Schwarz–Christoffel derivative and adding affine constants.

% BEGIN VOL04-EXPANSION IV18-example-02
\begin{example}[Right angles give square-root factors]\label{ex:iv18-expand-02}
For a rectangle, every finite vertex has interior angle \(\pi/2\), hence exponent
\[
\frac{\alpha}{\pi}-1=-\frac12.
\]
After normalizing three prevertices, the derivative therefore contains inverse square-root factors. Integrating leads to an elliptic integral, explaining why rectangle maps naturally connect Schwarz--Christoffel theory with elliptic functions.
\end{example}
% END VOL04-EXPANSION IV18-example-02
\section{Prevertices}
Determining the prevertices is the accessory-parameter problem and is generally nonlinear.

\section{Normalization}
Three real degrees of freedom can be fixed by half-plane Möbius automorphisms.

\section{Rectangles}
Symmetry reduces rectangle maps to elliptic integrals.

% BEGIN VOL04-EXPANSION IV18-example-03
\begin{example}[Triangle exponents]\label{ex:iv18-expand-03}
For a triangle with angles \(\alpha_1,\alpha_2,\alpha_3\), the derivative has factors \((z-x_k)^{\alpha_k-1}\). Since \(\alpha_1+\alpha_2+\alpha_3=\pi\), the three exponents sum to \(-2\). This global balance matches the behavior required at infinity for a half-plane-to-polygon map.
\end{example}
% END VOL04-EXPANSION IV18-example-03
\section{Slits and degenerate polygons}
Limiting polygon configurations produce maps to slit domains and strips.

\section{Numerical construction}
Practical Schwarz–Christoffel mapping often solves the parameter problem numerically.

% VOL04-NUMERICAL-SC-CONDITIONING-IV18
\begin{remark}[Accessory-parameter residuals and prevertex conditioning]
The Schwarz--Christoffel parameter problem is nonlinear.  Before solving it
numerically, the three real M\"obius degrees of freedom must be fixed by a
normalization; otherwise the parameterization contains a built-in
nonuniqueness and the Jacobian of the parameter equations can be singular for
purely representational reasons.

Let \(F(p)=0\) denote normalized side-length, cross-ratio, or other accessory
conditions for a parameter vector \(p\).  A computed parameter
\(\widetilde p\) should be accompanied by a residual
\[
\|F(\widetilde p)\|
\]
and by conditioning information for the Jacobian \(DF(\widetilde p)\), for
example its smallest singular value or condition number.  A small residual
does not imply small parameter error when the Jacobian is nearly singular.

Crowded prevertices are a characteristic source of sensitivity: widely
separated polygonal scales can correspond to prevertices that become extremely
close in the canonical domain.  In addition, the Schwarz--Christoffel integral
has algebraic endpoint behavior at prevertices, so generic quadrature without
branch-aware endpoint treatment can lose accuracy.

Numerical construction should therefore report the normalization, nonlinear
residual, stopping tolerance, Jacobian conditioning, ordering of prevertices,
branch convention, and quadrature strategy.  Geometric verification should
also check mapped vertex order and side lengths, not only the nonlinear
solver's residual.
\end{remark}

\section{Core structural results}

\begin{theorem}[Angle-exponent relation]\label{thm:iv18-01}
Near a prevertex $x_k$, a Schwarz–Christoffel map with derivative behaving like $(z-x_k)^{\alpha_k-1}$ maps a half-neighborhood to a wedge of interior angle $\alpha_k\pi$.
\begin{proof}
Integrating gives local behavior $F(z)-F(x_k)\sim C(z-x_k)^{\alpha_k}$. The argument is therefore multiplied by $\alpha_k$, sending the half-plane angle $\pi$ to $\alpha_k\pi$.
\end{proof}
\end{theorem}

\begin{theorem}[Sum of exponents]\label{thm:iv18-02}
For an $n$-gon, $\sum_{k=1}^n(\alpha_k-1)=-2$.
\begin{proof}
The interior-angle sum is $(n-2)\pi$, so $\sum\alpha_k=n-2$ and subtraction of $n$ gives $-2$.
\end{proof}
\end{theorem}

\begin{theorem}[Affine freedom]\label{thm:iv18-03}
Once the prevertices and exponents are fixed, two Schwarz–Christoffel maps with the same derivative factors differ by a nonzero multiplicative constant and an additive constant.
\begin{proof}
Their derivatives are proportional by the chosen normalization, so integration yields $F_2=A F_1+B$ with $A\ne0$.
\end{proof}
\end{theorem}

\section{Worked examples}

\begin{example}[Polygonal target diagnostic]\label{ex:iv18-worked-01}
Analyze Polygonal target in the setting of this chapter: state the precise hypothesis, translate polygon angles into Schwarz--Christoffel exponents and fix the affine constants from geometric normalization, and derive the resulting conclusion. A polygon is specified by ordered vertices and interior angles $\alpha_k\pi$. The chapter-specific mechanism is to translate polygon angles into Schwarz--Christoffel exponents and fix the affine constants from geometric normalization; this is the step that makes the conclusion genuinely complex-analytic.
\end{example}

\begin{example}[Derivative form diagnostic]\label{ex:iv18-worked-02}
Analyze Derivative form in the setting of this chapter: state the precise hypothesis, translate polygon angles into Schwarz--Christoffel exponents and fix the affine constants from geometric normalization, and derive the resulting conclusion. For a half-plane map with real prevertices $x_k$, the derivative has factors $(z-x_k)^{\alpha_k-1}$. The chapter-specific mechanism is to translate polygon angles into Schwarz--Christoffel exponents and fix the affine constants from geometric normalization; this is the step that makes the conclusion genuinely complex-analytic.
\end{example}

\begin{example}[Turning angles diagnostic]\label{ex:iv18-worked-03}
Analyze Turning angles in the setting of this chapter: state the precise hypothesis, translate polygon angles into Schwarz--Christoffel exponents and fix the affine constants from geometric normalization, and derive the resulting conclusion. The exponent $\alpha_k-1$ records the change of boundary direction at the corresponding vertex. The chapter-specific mechanism is to translate polygon angles into Schwarz--Christoffel exponents and fix the affine constants from geometric normalization; this is the step that makes the conclusion genuinely complex-analytic.
\end{example}

\begin{example}[Integration diagnostic]\label{ex:iv18-worked-04}
Analyze Integration in the setting of this chapter: state the precise hypothesis, translate polygon angles into Schwarz--Christoffel exponents and fix the affine constants from geometric normalization, and derive the resulting conclusion. The conformal map is obtained by integrating the Schwarz–Christoffel derivative and adding affine constants. The chapter-specific mechanism is to translate polygon angles into Schwarz--Christoffel exponents and fix the affine constants from geometric normalization; this is the step that makes the conclusion genuinely complex-analytic.
\end{example}

\section{Solved dossiers}
The dossiers are canonical solved problems. The provenance ledger separately records which retained corpus rules guided each topic and which dossiers were newly designed.

\begin{problem}[Polygonal target diagnostic]\label{prob:iv18-dossier-01}
Analyze Polygonal target in the setting of this chapter: state the precise hypothesis, translate polygon angles into Schwarz--Christoffel exponents and fix the affine constants from geometric normalization, and derive the resulting conclusion.
\end{problem}
\begin{solution}
A polygon is specified by ordered vertices and interior angles $\alpha_k\pi$. The chapter-specific mechanism is to translate polygon angles into Schwarz--Christoffel exponents and fix the affine constants from geometric normalization; this is the step that makes the conclusion genuinely complex-analytic.
\end{solution}

\begin{problem}[Derivative form diagnostic]\label{prob:iv18-dossier-02}
Analyze Derivative form in the setting of this chapter: state the precise hypothesis, translate polygon angles into Schwarz--Christoffel exponents and fix the affine constants from geometric normalization, and derive the resulting conclusion.
\end{problem}
\begin{solution}
For a half-plane map with real prevertices $x_k$, the derivative has factors $(z-x_k)^{\alpha_k-1}$. The chapter-specific mechanism is to translate polygon angles into Schwarz--Christoffel exponents and fix the affine constants from geometric normalization; this is the step that makes the conclusion genuinely complex-analytic.
\end{solution}

\begin{problem}[Turning angles diagnostic]\label{prob:iv18-dossier-03}
Analyze Turning angles in the setting of this chapter: state the precise hypothesis, translate polygon angles into Schwarz--Christoffel exponents and fix the affine constants from geometric normalization, and derive the resulting conclusion.
\end{problem}
\begin{solution}
The exponent $\alpha_k-1$ records the change of boundary direction at the corresponding vertex. The chapter-specific mechanism is to translate polygon angles into Schwarz--Christoffel exponents and fix the affine constants from geometric normalization; this is the step that makes the conclusion genuinely complex-analytic.
\end{solution}

\begin{problem}[Integration diagnostic]\label{prob:iv18-dossier-04}
Analyze Integration in the setting of this chapter: state the precise hypothesis, translate polygon angles into Schwarz--Christoffel exponents and fix the affine constants from geometric normalization, and derive the resulting conclusion.
\end{problem}
\begin{solution}
The conformal map is obtained by integrating the Schwarz–Christoffel derivative and adding affine constants. The chapter-specific mechanism is to translate polygon angles into Schwarz--Christoffel exponents and fix the affine constants from geometric normalization; this is the step that makes the conclusion genuinely complex-analytic.
\end{solution}

\begin{problem}[Prevertices diagnostic]\label{prob:iv18-dossier-05}
Analyze Prevertices in the setting of this chapter: state the precise hypothesis, translate polygon angles into Schwarz--Christoffel exponents and fix the affine constants from geometric normalization, and derive the resulting conclusion.
\end{problem}
\begin{solution}
Determining the prevertices is the accessory-parameter problem and is generally nonlinear. The chapter-specific mechanism is to translate polygon angles into Schwarz--Christoffel exponents and fix the affine constants from geometric normalization; this is the step that makes the conclusion genuinely complex-analytic.
\end{solution}

\begin{problem}[Normalization diagnostic]\label{prob:iv18-dossier-06}
Analyze Normalization in the setting of this chapter: state the precise hypothesis, translate polygon angles into Schwarz--Christoffel exponents and fix the affine constants from geometric normalization, and derive the resulting conclusion.
\end{problem}
\begin{solution}
Three real degrees of freedom can be fixed by half-plane Möbius automorphisms. The chapter-specific mechanism is to translate polygon angles into Schwarz--Christoffel exponents and fix the affine constants from geometric normalization; this is the step that makes the conclusion genuinely complex-analytic.
\end{solution}

\begin{problem}[Rectangles diagnostic]\label{prob:iv18-dossier-07}
Analyze Rectangles in the setting of this chapter: state the precise hypothesis, translate polygon angles into Schwarz--Christoffel exponents and fix the affine constants from geometric normalization, and derive the resulting conclusion.
\end{problem}
\begin{solution}
Symmetry reduces rectangle maps to elliptic integrals. The chapter-specific mechanism is to translate polygon angles into Schwarz--Christoffel exponents and fix the affine constants from geometric normalization; this is the step that makes the conclusion genuinely complex-analytic.
\end{solution}

\begin{problem}[Slits and degenerate polygons diagnostic]\label{prob:iv18-dossier-08}
Analyze Slits and degenerate polygons in the setting of this chapter: state the precise hypothesis, translate polygon angles into Schwarz--Christoffel exponents and fix the affine constants from geometric normalization, and derive the resulting conclusion.
\end{problem}
\begin{solution}
Limiting polygon configurations produce maps to slit domains and strips. The chapter-specific mechanism is to translate polygon angles into Schwarz--Christoffel exponents and fix the affine constants from geometric normalization; this is the step that makes the conclusion genuinely complex-analytic.
\end{solution}

\begin{problem}[Angle-exponent relation proof dossier]\label{prob:iv18-dossier-09}
Prove the following theorem-level statement and identify the essential hypothesis: Near a prevertex $x_k$, a Schwarz–Christoffel map with derivative behaving like $(z-x_k)^{\alpha_k-1}$ maps a half-neighborhood to a wedge of interior angle $\alpha_k\pi$.
\end{problem}
\begin{solution}
Integrating gives local behavior $F(z)-F(x_k)\sim C(z-x_k)^{\alpha_k}$. The argument is therefore multiplied by $\alpha_k$, sending the half-plane angle $\pi$ to $\alpha_k\pi$. The reusable mechanism here is to translate polygon angles into Schwarz--Christoffel exponents and fix the affine constants from geometric normalization.
\end{solution}

\begin{problem}[Sum of exponents proof dossier]\label{prob:iv18-dossier-10}
Prove the following theorem-level statement and identify the essential hypothesis: For an $n$-gon, $\sum_{k=1}^n(\alpha_k-1)=-2$.
\end{problem}
\begin{solution}
The interior-angle sum is $(n-2)\pi$, so $\sum\alpha_k=n-2$ and subtraction of $n$ gives $-2$. The reusable mechanism here is to translate polygon angles into Schwarz--Christoffel exponents and fix the affine constants from geometric normalization.
\end{solution}

\begin{problem}[Affine freedom proof dossier]\label{prob:iv18-dossier-11}
Prove the following theorem-level statement and identify the essential hypothesis: Once the prevertices and exponents are fixed, two Schwarz–Christoffel maps with the same derivative factors differ by a nonzero multiplicative constant and an additive constant.
\end{problem}
\begin{solution}
Their derivatives are proportional by the chosen normalization, so integration yields $F_2=A F_1+B$ with $A\ne0$. The reusable mechanism here is to translate polygon angles into Schwarz--Christoffel exponents and fix the affine constants from geometric normalization.
\end{solution}

\begin{problem}[Chapter synthesis]\label{prob:iv18-dossier-12}
Connect the definitions and principal theorems of this chapter into one reusable complex-analysis strategy.
\end{problem}
\begin{solution}
Schwarz–Christoffel transformations map canonical half-plane or disk domains to polygonal regions. The exponents in the derivative encode turning angles at polygon vertices. A chapter-specific strategy is to translate polygon angles into Schwarz--Christoffel exponents and fix the affine constants from geometric normalization; the payoff is that the derivative singularities encode the turning angles of the target polygon.
\end{solution}

\section{Exercises with complete solutions}

\begin{exercise}\label{exr:iv18-01}
Use Polygonal target to carry out the chapter's main method: translate polygon angles into Schwarz--Christoffel exponents and fix the affine constants from geometric normalization. State one conclusion whose validity depends on the relevant hypotheses.
\end{exercise}
\begin{hint}
For a polygon vertex with interior angle \(\alpha_k\pi\), the Schwarz--Christoffel derivative contributes exponent \(\alpha_k-1\).
\end{hint}
\begin{solution}
A polygon is specified by ordered vertices and interior angles $\alpha_k\pi$. In this chapter the concrete payoff is that the derivative singularities encode the turning angles of the target polygon.
\end{solution}

\begin{exercise}\label{exr:iv18-02}
Use Derivative form to carry out the chapter's main method: translate polygon angles into Schwarz--Christoffel exponents and fix the affine constants from geometric normalization. State one conclusion whose validity depends on the relevant hypotheses.
\end{exercise}
\begin{hint}
Place prevertices in boundary order before integrating the product formula; wrong ordering changes the polygon.
\end{hint}
\begin{solution}
For a half-plane map with real prevertices $x_k$, the derivative has factors $(z-x_k)^{\alpha_k-1}$. In this chapter the concrete payoff is that the derivative singularities encode the turning angles of the target polygon.
\end{solution}

\begin{exercise}\label{exr:iv18-03}
Use Turning angles to carry out the chapter's main method: translate polygon angles into Schwarz--Christoffel exponents and fix the affine constants from geometric normalization. State one conclusion whose validity depends on the relevant hypotheses.
\end{exercise}
\begin{hint}
Use Möbius freedom to fix three prevertices and reduce the number of accessory parameters.
\end{hint}
\begin{solution}
The exponent $\alpha_k-1$ records the change of boundary direction at the corresponding vertex. In this chapter the concrete payoff is that the derivative singularities encode the turning angles of the target polygon.
\end{solution}

\begin{exercise}\label{exr:iv18-04}
Use Integration to carry out the chapter's main method: translate polygon angles into Schwarz--Christoffel exponents and fix the affine constants from geometric normalization. State one conclusion whose validity depends on the relevant hypotheses.
\end{exercise}
\begin{hint}
Near a prevertex, integrate \((z-x_k)^{\alpha_k-1}\) to recover the local corner angle.
\end{hint}
\begin{solution}
The conformal map is obtained by integrating the Schwarz–Christoffel derivative and adding affine constants. In this chapter the concrete payoff is that the derivative singularities encode the turning angles of the target polygon.
\end{solution}

\begin{exercise}\label{exr:iv18-05}
Use Prevertices to carry out the chapter's main method: translate polygon angles into Schwarz--Christoffel exponents and fix the affine constants from geometric normalization. State one conclusion whose validity depends on the relevant hypotheses.
\end{exercise}
\begin{hint}
Determine multiplicative and additive constants from side lengths, orientation, and one base point.
\end{hint}
\begin{solution}
Determining the prevertices is the accessory-parameter problem and is generally nonlinear. In this chapter the concrete payoff is that the derivative singularities encode the turning angles of the target polygon.
\end{solution}

\begin{exercise}\label{exr:iv18-06}
Use Normalization to carry out the chapter's main method: translate polygon angles into Schwarz--Christoffel exponents and fix the affine constants from geometric normalization. State one conclusion whose validity depends on the relevant hypotheses.
\end{exercise}
\begin{hint}
For unbounded polygons, identify which prevertex corresponds to infinity and adjust the exponent bookkeeping.
\end{hint}
\begin{solution}
Three real degrees of freedom can be fixed by half-plane Möbius automorphisms. In this chapter the concrete payoff is that the derivative singularities encode the turning angles of the target polygon.
\end{solution}

\begin{exercise}\label{exr:iv18-07}
Use Rectangles to carry out the chapter's main method: translate polygon angles into Schwarz--Christoffel exponents and fix the affine constants from geometric normalization. State one conclusion whose validity depends on the relevant hypotheses.
\end{exercise}
\begin{hint}
Check branch choices of the fractional powers along the boundary so side directions are consistent.
\end{hint}
\begin{solution}
Symmetry reduces rectangle maps to elliptic integrals. In this chapter the concrete payoff is that the derivative singularities encode the turning angles of the target polygon.
\end{solution}

\begin{exercise}\label{exr:iv18-08}
Use Slits and degenerate polygons to carry out the chapter's main method: translate polygon angles into Schwarz--Christoffel exponents and fix the affine constants from geometric normalization. State one conclusion whose validity depends on the relevant hypotheses.
\end{exercise}
\begin{hint}
Differentiate a candidate map first; verifying the Schwarz--Christoffel derivative is often simpler than integrating explicitly.
\end{hint}
\begin{solution}
Limiting polygon configurations produce maps to slit domains and strips. In this chapter the concrete payoff is that the derivative singularities encode the turning angles of the target polygon.
\end{solution}

\begin{exercise}[Small residual with an ill-conditioned accessory system]\label{exr:iv18-numerical-sc-01}
A normalized Schwarz--Christoffel parameter solve has residual
\[
\|F(\widetilde p)\|_2\le10^{-8}.
\]
At the computed point, suppose the square Jacobian is invertible with
\[
\sigma_{\min}(DF(\widetilde p))=10^{-5}.
\]
Using the linearized correction
\[
DF(\widetilde p)\,\Delta p=-F(\widetilde p),
\]
bound \(\|\Delta p\|_2\).  What does this calculation show about interpreting a
small nonlinear residual?
\end{exercise}
\begin{hint}
For an invertible matrix \(J\),
\[
\|J^{-1}\|_2=1/\sigma_{\min}(J).
\]
Apply this to the linearized correction equation.
\end{hint}
\begin{solution}
Here
\[
\|DF(\widetilde p)^{-1}\|_2=10^5,
\]
so the linearized correction satisfies
\[
\|\Delta p\|_2
\le
10^5\,10^{-8}
=
10^{-3}.
\]
Thus a residual of size \(10^{-8}\) can still permit a much larger
first-order parameter correction when the accessory equations are poorly
conditioned.  The estimate concerns the local linearized correction; a
rigorous nonlinear parameter-error bound would additionally require control of
Jacobian variation and validity of the local inverse model.
\end{solution}

\section*{Chapter summary}
The chapter now forms part of the canonical complex-analysis chain from holomorphicity through residues, global continuation, and Riemann surfaces.
% BEGIN VOL04-EXPANSION IV18-exercises-01
\section{Graded supplementary exercises}
These exercises extend the protected chapter with a balanced set of computations, proofs, hypothesis tests, applications, and challenges.

\subsection*{Standard computations}

\begin{exercise}[Right angle exponent]\label{exr:iv18-expand-01}
Find the Schwarz--Christoffel exponent for an interior angle \(\pi/2\).
\end{exercise}
\begin{hint}
Compute \(\alpha/\pi-1\).
\end{hint}
\begin{solution}
The exponent is \(-1/2\).
\end{solution}

\begin{exercise}[Equilateral triangle]\label{exr:iv18-expand-02}
Find the exponent at each vertex of an equilateral triangle.
\end{exercise}
\begin{hint}
Each interior angle is \(\pi/3\).
\end{hint}
\begin{solution}
Each exponent is \(-2/3\).
\end{solution}

\begin{exercise}[Reflex vertex]\label{exr:iv18-expand-03}
Find the exponent for interior angle \(3\pi/2\).
\end{exercise}
\begin{hint}
Use the same formula.
\end{hint}
\begin{solution}
The exponent is \(1/2\), reflecting a reflex corner.
\end{solution}

\begin{exercise}[Straight angle]\label{exr:iv18-expand-04}
What exponent corresponds to interior angle \(\pi\)?
\end{exercise}
\begin{hint}
Substitute directly.
\end{hint}
\begin{solution}
The exponent is zero, so no singular power factor is needed for a straight continuation.
\end{solution}

\begin{exercise}[Triangle sum]\label{exr:iv18-expand-05}
What is the sum of the three derivative exponents for a triangle?
\end{exercise}
\begin{hint}
Use the angle sum \(\pi\).
\end{hint}
\begin{solution}
The sum is \((\pi/\pi)-3=-2\).
\end{solution}

\subsection*{Proofs}

\begin{exercise}[Local corner law]\label{exr:iv18-expand-06}
Derive the local angle law from \(f'(z)\sim C(z-x_k)^{\beta_k-1}\).
\end{exercise}
\begin{hint}
Integrate the leading power.
\end{hint}
\begin{solution}
Integration gives \(f(z)-f(x_k)\sim (C/\beta_k)(z-x_k)^{\beta_k}\). Arguments are multiplied by \(\beta_k\), so the half-plane angle \(\pi\) becomes \(\beta_k\pi=\alpha_k\).
\end{solution}

\begin{exercise}[Polygon exponent sum]\label{exr:iv18-expand-07}
For an \(n\)-gon, show that the finite-vertex exponents \(\alpha_k-1\) sum to \(-2\).
\end{exercise}
\begin{hint}
Use the polygon interior angle sum.
\end{hint}
\begin{solution}
Since \(\sum\alpha_k=(n-2)\pi\), the exponent sum is \((n-2)-n=-2\).
\end{solution}

\begin{exercise}[Affine freedom]\label{exr:iv18-expand-08}
Explain why multiplying a Schwarz--Christoffel map by a nonzero constant and adding a constant preserves the polygon shape up to similarity and translation.
\end{exercise}
\begin{hint}
Differentiate the transformed map.
\end{hint}
\begin{solution}
The derivative is multiplied by one nonzero complex constant, which rotates and scales all edges uniformly; the additive constant translates the image.
\end{solution}

\begin{exercise}[Prevertex order]\label{exr:iv18-expand-09}
Why must the real prevertices occur in the same cyclic order as the polygon vertices for a standard upper-half-plane map?
\end{exercise}
\begin{hint}
Follow the boundary image of the oriented real axis.
\end{hint}
\begin{solution}
The real boundary is traversed monotonically through the prevertices. Continuity maps successive real intervals to successive polygon edges, fixing the cyclic vertex order.
\end{solution}

\subsection*{Counterexamples and hypothesis tests}

\begin{exercise}[Wrong exponent]\label{exr:iv18-expand-10}
What geometric defect results if a right-angle vertex is assigned exponent \(-1/3\) instead of \(-1/2\)?
\end{exercise}
\begin{hint}
Convert the exponent back to an angle.
\end{hint}
\begin{solution}
The corresponding angle factor is \(2/3\), giving an interior angle \(2\pi/3\) rather than \(\pi/2\).
\end{solution}

\begin{exercise}[Repeated prevertices]\label{exr:iv18-expand-11}
Can two distinct finite polygon vertices be assigned the same ordinary prevertex in a nondegenerate Schwarz--Christoffel map?
\end{exercise}
\begin{hint}
A single boundary point has only one local turning exponent.
\end{hint}
\begin{solution}
Not in the standard nondegenerate setup. Coincident prevertices signal a limiting or degenerate polygon configuration.
\end{solution}

\begin{exercise}[Ignoring normalization]\label{exr:iv18-expand-12}
Why are prevertex positions not all independent parameters?
\end{exercise}
\begin{hint}
The upper half-plane has a three-real-parameter Möbius automorphism group.
\end{hint}
\begin{solution}
Three real degrees of freedom can be normalized, commonly by fixing three prevertices. Failing to account for this produces redundant parameters.
\end{solution}

\subsection*{Applications and investigations}

\begin{exercise}[Numerical parameter problem]\label{exr:iv18-expand-13}
What must be solved numerically when mapping the half-plane to a specified polygon with more than three vertices?
\end{exercise}
\begin{hint}
Angles determine exponents, but side lengths determine prevertex spacing.
\end{hint}
\begin{solution}
One solves nonlinear equations for the remaining prevertices and scale so that integrals between consecutive prevertices reproduce the prescribed side lengths.
\end{solution}

\begin{exercise}[Flow around polygons]\label{exr:iv18-expand-14}
Why are Schwarz--Christoffel maps useful in two dimensional potential flow?
\end{exercise}
\begin{hint}
Conformal maps preserve harmonic structure and angles.
\end{hint}
\begin{solution}
A difficult polygonal boundary can be mapped to a half-plane where analytic potentials are simple, then transported back to the physical polygonal domain.
\end{solution}

\subsection*{Challenge problems}

\begin{exercise}[Rectangle and elliptic integrals]\label{exr:iv18-expand-15}
Explain why the rectangle map leads to an integral with four square-root branch points.
\end{exercise}
\begin{hint}
Each right angle contributes exponent \(-1/2\).
\end{hint}
\begin{solution}
With four finite prevertices in a symmetric normalization, the derivative is proportional to the reciprocal square root of a quartic polynomial. Integrating this differential is an elliptic integral.
\end{solution}

\begin{exercise}[Vertex at infinity]\label{exr:iv18-expand-16}
How can placing one polygon vertex at infinity simplify a Schwarz--Christoffel formula?
\end{exercise}
\begin{hint}
Use Möbius freedom to send a convenient prevertex to infinity.
\end{hint}
\begin{solution}
Its factor is absorbed into the behavior at infinity, reducing the number of explicit finite factors and often simplifying both the integral and the parameter equations.
\end{solution}

% END VOL04-EXPANSION IV18-exercises-01
